\chapter{Planteamiento del Problema y Justificación}
\begin{sloppypar}
    \section{Planteamiento del problema}
    Dado el modelo físico del sistem de levitación magnética en configuración inestable con efecto atascamiento-deslizamiento 
    diseñar un controlador de posición bajo el enfoque de control lineal.
    Se aplicarán técnicas de control clásico para proponer un controlador con doble efecto integral (\textbf{PII}). 
    
    \section{Justificación}
    Controlar sistemas de levitación magnética con componentes mecánicos implica que se conviertan en complejos sistemas de ingeniería en los que la comunidad científica se ha interesado por años.
    La zona muerta, como miembro de la famila de las no linealidades no suaves, impone un reto
    al uso de controladores lineales para configuraciones de sistemas de este estilo. El mencionado fenómeno es capaz de causar inestabilidad en el sistema, 
    deterioro del desempeño en estado estable en lazo cerrado especialmente al alrededor del punto de equilibrio donde los desplazamientos suelen ser muy pequeños
    y por tanto se opera dentro de esta zona fuerza al sistema a presentar oscilaciones de alta 
    frecuencia en la variable de control y si se adiciona la dinámica no lineal intrínsica de los electroimanes,
    se llegan a escenarios donde el diseño del esquema de control se convierte en un reto. Trabajos como \cite{khimani},\cite{mukesh},\cite{neelma}, evitan incluir en su modelado
    la zona muerta ya sea porque prefieren ignorarla o incluirlas como dinámicas no modelas y/o incertidumbres no estructuradas. Las estrategias de control que se han propuesto para resolver 
    esta problemática van desde proponer controladores robustos adaptativos difusos, QFT hasta PID \cite{kuo1},\cite{mukesh},\cite{maaz} que tienen un buen manejo del error en estado estable
    pero aun así la literatura es escasa. En lo que respecta a los modelos que se han propuesto para representar la zona muerta se encuentran los modelos lineales \cite{salim}, no lineales \cite{huanqing}, difusos \cite{zhi},
    asimétricos \cite{gang}, simétricos \cite{nizar}, generalizados \cite{guanyu}, entre otros; desafortunadamente estos modelos y junto con los resultados que se muestran en los experimentos no son 
    completamente representativos de la realidad. 
    
    Se propone el diseño de un lazo de control con técnicas de control clásico que satisfaga condiciones de seguimiento en un sistema de levitación magnética. Alcanzar un comportamiento robusto
    en el sistema de control en donde el diseño del controlador de tipo \textbf{P}roporcional \textbf{I}ntegral \textbf{I}ntegral (\textbf{PII}) considere la fricción y la zona muerta dinámica
    la cual depende de la posición donde se encuentra el cuerpo levitante, es el reto propuesto.
\end{sloppypar}